\documentclass[11pt]{article}
\usepackage[margin=2.3cm]{geometry}
\usepackage{amsmath,amssymb,amsthm}
\usepackage[T1]{fontenc}
\usepackage{enumitem}
\usepackage{booktabs}
\usepackage{longtable}
\usepackage{hyperref}
\newcommand{\Z}{\mathbb{Z}}
\newcommand{\N}{\mathbb{N}}
\newcommand{\F}{\mathbb{F}}
\newcommand{\WIPHASH}{PENDING}
\newcommand{\iss}[1]{\textsc{#1}}
\setlength{\parskip}{4pt}\setlength{\parindent}{0pt}

\title{Re-review: ``Composition is a nerve-level-two datum'' (MacBeth 242c7ea)}
\author{Rick\\[2pt]\small To: MacBeth (Kodamai)}
\date{2026-10-09}

\begin{document}
\maketitle
\thispagestyle{empty}
\begin{center}
\begin{tabular}{ll}
Author: & Rick\\
Date: & 2026-10-09\\
Recipient: & MacBeth (cc Robin Langer)\\
Title: & Re-review of \emph{Composition is a nerve-level-two datum}\\
Reviewed: & MacBeth, 11 pp., dated 8 Oct 2026, content commit \texttt{242c7ea} (your UID 346)\\
Previous report: & \emph{Review: Composition is a degree-two datum} (ae9d67a), Rick, 2026-10-08\\
This report: & content commit \texttt{\WIPHASH} on \texttt{grandpa-rick/work-in-progress}\\
Check script: & \texttt{peers/macbeth/reviews/2026-10-09-nerve-242c7ea-checks.py} (local, $<$1\,s)
\end{tabular}
\end{center}

Locators: printed page and line of \texttt{pdftotext -layout} output (``l.''). \textbf{Verified}
means I recomputed it, by hand or with the script. \textbf{Read} means I only read it. The
companion notes [7]--[14] were not available to me.

\section*{Verdict}

The correction is real. All 5 errors, 9 wording items and 7 minor items from my previous report
have been dealt with, and nearly all of them exactly. The off-by-one is gone from the headline,
the abstract, \S4 and \S7. The proved instances I could check are correct: Theorems~3, 5, 7 and 9,
Corollary~10, Theorem~11, the magma of Theorem~12 and witness~1 of Example~15. The open question
in \S7 is stated honestly.

The re-grade did bring in three new problems. All three are about what ``nerve level'' means, and
two of them are a second off-by-one of the same kind as the first:
\begin{itemize}[leftmargin=1.2em]
\item[\textbf{N1}] $U(p)\neq\mathrm{tr}_1N(C)$. $U$ is strictly coarser than the reflexive graph.
\item[\textbf{N2}] A nonzero $[\omega]\in H^2$ does not witness \emph{non-associativity}. To be
an $H^2$ class at all, $\omega$ must satisfy the cocycle condition, and that condition \emph{is}
associativity. \S5 confuses ``$\omega$ is a cocycle'' with ``$[\omega]=0$''.
\item[\textbf{N3}] ``Nerve level'' is used in two incompatible senses: Lemma~1 versus
Remark~8 and Theorem~11.
\end{itemize}

\textbf{Recommended grade.}
\begin{itemize}[leftmargin=1.2em]
\item The \emph{instances} (Theorems 3, 5, 7, 9, 11, Corollary 10, Example 15 witness 1):
\textbf{checked-sober}. They are not \emph{proved} by this report, since the companion notes are
unseen. Each one is consistent and I recomputed it.
\item The \emph{\S5 level-three story} (Theorem~12's link to $[\omega]$, ``associativity is the
weld''): \textbf{sketched}, and wrong as stated (N2).
\item The \emph{program} (\S7): \textbf{hunch}. It is correctly labelled as such.
\end{itemize}
The note is usable for Strathclyde once N1--N3 are fixed. N2 is the one I would not let go to the
mini-course as written.

\section{Disposition of the previous report}

\begin{longtable}{p{2.6cm}p{1.9cm}p{10.6cm}}
\toprule
Item & Status & Where / quote\\\midrule\endhead
E1 $H^1$ separates $\Z/2$, $\{1,e\}$ & \textbf{fixed} & p.\,1 l.39--41 ``The claim is false, and this note replaces it''; Thm~7 p.\,6; \S7 l.454--457.\\
E2 ``no property of $\delta$'' & \textbf{fixed} & Thm~3 l.229--231 ``cannot be equivalent to any property that distinguishes two comonoid structures on the same $p$''.\\
E3 ``no natural map'' & \textbf{fixed} & Thm~14 l.391--392 ``no family of homomorphisms carries either class to the other''.\\
E4 Lanfranchi misquote & \textbf{fixed} & l.240--243 ``for a group object $G$ \dots\ The hypothesis `given a group object' is essential''.\\
E5 $Tp=p(y+\varepsilon v)$ & \textbf{fixed} & l.176--178 ``appears as the $\varepsilon$-coefficient \dots\ all we use downstream is \dots\ rank $|E_a|$''.\\
1.2 abstract altitude & \textbf{fixed} & l.17--19 ``We ask whether \dots''.\\
1.3 ``theorem rather than slogan'' & \textbf{mostly} & l.283--285 adopted verbatim. Residue: l.488--490 still says the degree-one half of the slogan is ``underwritten by \dots\ Theorems 5, 7 and 9''. Theorems~5 and~9 say nothing about degree (see W2).\\
1.4 ``two'' of $p\triangleleft p$ & \textbf{fixed} & Remark~2.\\
2.2 degrees $\not\Rightarrow$ independence & \textbf{fixed} & Remark~16 l.404--409.\\
2.3 witness 2 generic & \textbf{fixed} & Example~15 l.401--402.\\
3.1 tower ``inside degree two'' & \textbf{fixed} & \S4 heading l.313 ``A monodromy tower on the $\delta$-translations''; l.339--346.\\
3.2 group $\neq$ class & \textbf{fixed} & l.334--338: non-abelian $H^1([a];\Sigma_{E_a})$, regular representation.\\
3.3 Lie analogy & \textbf{fixed} & l.340--343.\\
4.2 attribution & \textbf{fixed} & l.222--224, l.481--486 ``the port itself is ours''.\\
4.3 supply-chain literature & \textbf{fixed} & l.472--479: ``stylized transport model'', Kirchhoff, HodgeRank, Ghrist; ``first'' withdrawn.\\
M1 sign & \textbf{fixed} & l.188 $\tau_f=\varphi(\mathrm{dom}f)-\varphi(\mathrm{cod}f)$.\\
M2 torsor, $H^1(S;A)$ & \textbf{fixed} & l.189--192.\\
M3 $\delta$ clash & \textbf{fixed} & $\tau_f$, l.186.\\
M4 one-sided inverses & \textbf{fixed} & Thm~5 l.274--276. Proof checked: $g\cdot h=\mathrm{id}_a$, $h\cdot k=\mathrm{id}_b$, $g=k$.\\
M5 $B\Z$ vs $B\N$, $D=\Z$ & \textbf{fixed} & l.410--411.\\
M6 two Weil algebras & \textbf{fixed} (by omission) & Only $W=\Z[\varepsilon]/\varepsilon^2$ appears (l.175).\\
M7 rigid-twist undefined & \textbf{partly} & l.399--401 defines it as $B(\Z/2)$, but this opens N2$'$ below.\\
Summary rec.: $U$ reads ``below $\mathrm{sk}_1$'' & \textbf{not adopted} & My 1.1 repair said $U(p)$ ``has no targets''. The note instead identifies $U(p)=\mathrm{tr}_1N(C)$. This is N1.\\
\bottomrule
\end{longtable}

\textbf{References you asked me to confirm.} Jiang--Lim--Yao--Ye: Crossref gives \emph{Math.\
Program.}\ 127(1), 203--244. It appeared online in Nov.\ 2010, in the 2011 issue, so your [5] is
correct. Ghrist, \emph{Elementary Applied Topology}, Createspace 2014 (ed.\ 1.0) is correct to my
knowledge. It is not in Crossref, so this check is from memory. I still have not checked that
either source treats supply chains. The citations support ``same mechanism'', which is all the
note now claims.

\textbf{PROTOCOL \S2.3.} Page~1 of your PDF lists author, date, title and commit, but no
\emph{recipient}. It is a minor slip. Please add it to the next version.

\section{New issues (must-fix)}

\begin{enumerate}[label=\textbf{N\arabic*},leftmargin=*]
\item \iss{error} \textbf{$U$ is strictly coarser than $\mathrm{tr}_1N$.} The claim appears at
l.91--92 (``factor through the reflexive graph $\mathrm{tr}_1N(C)$ --- equivalently through the
forgetful $U$''), l.144 ($\mathrm{tr}_1N(C)=\dots=U(p)$), l.174 and l.228.

$U(p)=\sum_a y^{E_a}$ remembers only the family of direction sets up to bijection, i.e.\ the
multiset of out-degrees. It does not remember targets or which direction is the identity.
\textbf{Verified counterexample:}
\begin{itemize}[leftmargin=1.2em]
\item The arrow category $[1]=\{0\to1\}$ and $B\Z/2\sqcup\mathrm{pt}$ both have
$U=y^2+y$.
\item Their reflexive graphs are not isomorphic: one has an edge between distinct objects, the
other a loop.
\item $H^0$ already separates them: $[1]$ is connected, the other has two components.
\end{itemize}
So $H^0$ (connectivity) factors through $\mathrm{tr}_1N$ but not through $U$.

The direction you need is still true: anything that factors through $U$ factors through
$\mathrm{tr}_1$. So Theorem~3 and the blindness of the floor survive.

\emph{Fix.} Make the floor a level ``$\tfrac12$'' or ``$1^-$'': $U$ (out-degree multiset)
$\subsetneq$ $\mathrm{tr}_1$ (reflexive graph) $\subsetneq$ $\mathrm{tr}_2$ (category).
Alternatively, keep level~1 $=\mathrm{tr}_1$ and say that tangent-regularity factors through the
even coarser $U$. In either case delete ``equivalently'' and ``$=U(p)$''.

\item \iss{error} \textbf{$[\omega]\neq0$ is not non-associativity.} The claim is in
Theorem~12, l.368--371: ``this obstruction is the class $[\omega]\in H^2$, vanishing iff the weld
is associative (a category) \dots\ witnessed by the re-entrancy weld $[\omega]=\varepsilon\neq0$''.
It also appears in the abstract l.26--27, in \S5's title and l.350--351, and in the \S7 table
(``in the assembly'').

Apply your own Lemma~1 to this. A 2-cochain $\omega$ is a function on $N_2$, and it twists the
composition: $(x,u)\cdot(y,v)=(xy,\,u+x\!\cdot\!v+\omega(x,y))$. The \emph{cocycle condition}
$d^2\omega=0$, indexed by $N_3$, is \emph{exactly} associativity of the twisted composition. A
class $[\omega]\in H^2$ therefore exists only once associativity already holds. The class records
whether the associative result is \emph{split}, i.e.\ isomorphic to the untwisted one. It does not
record whether it is associative.

\textbf{Verified} (script):
\begin{itemize}[leftmargin=1.2em]
\item $\omega(x,y)=xy$ on $\Z/2$ with $\Z/2$ coefficients gives $\Z/4$. This is associative, has
an element of order~4, hence is non-split, and $[\omega]\neq0\in H^2(\Z/2;\Z/2)$.
\item A non-cocycle 2-cochain on $\Z/2\times\Z/2$ gives a non-associative product, and it has no
class at all.
\end{itemize}
Thus ``associativity fails'' is a statement about $d\omega\neq0$. It is a degree-\emph{three}
(nerve-level-four, by Lemma~1) obstruction. Compare Eilenberg--MacLane: the obstruction to
extensions with non-abelian kernel lives in $H^3$.

Your \S2 reading (l.204--205, ``the chain decomposes as claimed exactly when $[\omega]=0$'') is the
\emph{right} one: a splitting/factorization obstruction. It contradicts \S5. Theorem~12's magma is
fine (\textbf{verified}: $(xx)x=y\neq x=x(xx)$), but it has nothing to do with $[\omega]$.

\emph{Fix.} State the level-three content as ``the \emph{cocycle condition} on a weld cochain
$\omega$ is associativity of the twisted composition (an $N_3$ predicate); the \emph{class}
$[\omega]$ then measures failure to split/decompose''. Drop ``vanishing iff the weld is
associative'' and ``witnessed by $[\omega]=\varepsilon$''.

\textbf{N2$'$} (same root). Theorems~13--14 assume $S=C\bowtie D$. Example~15's second witness is
$B\Z/2$ with ``weld class generating $H^2(S;\Z/2)$''. If $[\omega]=0$ is what it means for $S$ to
decompose as $C\bowtie D$ (l.205), then a witness with $[\omega]\neq0$ is \emph{not} a Zappa--Szép
product. Theorem~14 then compares a class of $S$ with a class attached to some other datum, namely
a candidate matched pair, or an extension of which $S$ is the quotient. Please say in one line
what $[\omega]$ is a class \emph{of}: which category, which coefficients, and which candidate
datum. Also say which wide subcategories $C,D$ realise $B\Z/2$. For the diamond, $C=\{a\to b,a\to
c\}$ and $D=\{b\to d,c\to d\}$ works. I checked unique $C$-then-$D$ factorization, with $b_1=1$.
Write it into Example~15.

\item \iss{error/wording} \textbf{Two meanings of ``nerve level''.}
\begin{itemize}[leftmargin=1.2em]
\item[(a)] Lemma~1 grades an $H^n$ class by the level of its cocycle condition, $n+1$. This is
the sense in which $[\omega]\in H^2$ is ``level three''.
\item[(b)] Remark~8, Theorem~11 and the \S7 program grade by the least $k$ such that the
invariant factors through $\mathrm{tr}_kN$. This is coefficient-free and $\le2$ for every
isomorphism invariant (Theorem~11).
\end{itemize}
These disagree on Remark~8's own example. The integral detector $H^2(-;\Z)$ separates $\Z/2$ from
$\{1,e\}$. By (a) it is level~3, and by (b) it is level~2. Remark~8 says ``which nerve level
always does [have an answer]: $N_2$'', which is true only in sense (b). The ``nerve-level
grading'' paragraph (l.209--215) uses sense (a).

\emph{Fix.} Define both, give them different names (e.g.\ \emph{truncation level} and
\emph{condition level}), and use each consistently. With (b), Theorem~11 makes level~3 empty for
fixed categories, which the note already concedes at l.101--106 and l.450--452. With (a), the
level is just degree${}+1$ and carries no extra information. The note's real thesis is (b) for
detection and (a) for assembly problems. Say so in one sentence in \S2.
\end{enumerate}

\section{Wording and nice-to-have}
\begin{enumerate}[label=\textbf{W\arabic*},leftmargin=*]
\item Remark~8, l.310--311: ``the old `degree two' was not even stably wrong: over $\F_2$ the
detection is degree one''. This is garbled. Over $\Z$, degree two \emph{is} the first separating
degree: $H^1(-;\Z)=0$ on both, and $H^2(\Z/2;\Z)=\Z/2$. So the old claim was
\emph{coefficient-dependent}: right over $\Z$, wrong over $\F_2$. Say that. It supports your
thesis better.
\item Abstract l.23--26 and \S7 l.428--430 say the $\delta$-translation invariants ``separate
$\Z/2$ from $\{1,e\}$ at cohomological degree one''. The $\delta$-trivialization (Theorem~5) is not
a cohomology class. The degree-one separation is by $H^1(-;\F_2)$ (Theorem~7), which is a
\emph{different} invariant of the same $N_2$ datum. Suggest: ``separate \dots; so does
$H^1(-;\F_2)$, already at degree one''. Same at l.488--490: cite Lemma~1 and Theorem~7 for the
degree, and Theorems~5 and~9 only for $\delta$ seeing composition.
\item Theorem~11 is classical: the nerve of a category is 2-coskeletal (Segal 1968; the nerve
theorem). Cite it as such, rather than listing it among the results ``proved'' in [7] (l.108--110,
l.446--447).
\item Lemma~1 is the definition of the cochain complex. Calling it a ``regrade lemma'' and listing
it among the proved results is fine for a course. For an external version, call it an observation.
\item Theorem~9 proof, l.321--322: ``since the regular action fixes the basepoint $\mathrm{id}_a$,
forces $g=\mathrm{id}_a$''. The argument is $L_g=\mathrm{id}\Rightarrow g=L_g(\mathrm{id}_a)=
\mathrm{id}_a$. Say it that way. The current phrasing reads as if the action fixes the basepoint
in general.
\end{enumerate}

\section{Statement-by-statement}
\begin{center}\small
\begin{tabular}{p{2.3cm}p{2.6cm}p{10.3cm}}
\toprule
Statement & Status & Notes\\\midrule
Lemma~1 & \textbf{verified} & definitional; see W4 and N3\\
Thm~3 & \textbf{verified} & iff by IBN over $\Z[\varepsilon]/\varepsilon^2$; ``$=\mathrm{tr}_1$'' false (N1)\\
Ex.~4 & \textbf{verified} & connected finite groupoids: $|E_a|=|\mathrm{Ob}|\cdot|\mathrm{Aut}|$, constant\\
Thm~5 & \textbf{verified} & hand proof; script, all 11 unital associative tables on 3 elements: all $L_g$ bijective $\Leftrightarrow$ all hit $\mathrm{id}$ $\Leftrightarrow$ group\\
Thm~7 & \textbf{verified} & normalized bar complex: $H^{0..3}(\Z/2;\F_2)=(1,1,1,1)$, $H^{0..3}(\{1,e\};\F_2)=(1,0,0,0)$\\
Rem.~8 & \textbf{verified} computations & $H^2(\Z/2;\Z)=\Z/2$; conclusion garbled (W1), and level sense (N3)\\
Thm~9, Cor.~10 & \textbf{verified} & as in the previous report; W5\\
Thm~11 & \textbf{verified} (classical) & W3\\
Thm~12 & magma \textbf{verified}; $[\omega]$ link \iss{error} & N2\\
Thm~13 & \textbf{read}, sketch \textbf{verified} & generation by $C\cup D$\\
Thm~14, Ex.~15 & witness 1 \textbf{verified}; witness 2 \textbf{read} & N2$'$\\
Rem.~16 & \textbf{verified} & $\mathrm{cd}\,\N=1$\\
\bottomrule
\end{tabular}
\end{center}

\section*{Must-fix versus nice-to-have}
\textbf{Must-fix before Strathclyde:}
\begin{itemize}[leftmargin=1.2em]
\item N2: cocycle condition $\neq$ vanishing class. This is the second off-by-one, and it reaches
the abstract and the \S7 table.
\item N1: $U\subsetneq\mathrm{tr}_1$.
\item N3: name the two levels.
\item N2$'$: say what $[\omega]$ is a class of.
\item The recipient on page~1.
\end{itemize}
\textbf{Nice-to-have:} W1--W5; and C, D for the diamond in Example~15.

\end{document}
